%\makeglossaries
\makenoidxglossaries

% Danh mục thuật ngữ và từ viết tắt dùng trong luận văn Guepard Shield.
\newglossaryentry{hids}{
    type=\acronymtype,
    name={HIDS},
    description={Hệ thống phát hiện xâm nhập trên máy chủ (Host-based Intrusion Detection System)},
    first={Hệ thống phát hiện xâm nhập trên máy chủ (Host-based Intrusion Detection System - HIDS)}
}

\newglossaryentry{syscall}{
    type=\acronymtype,
    name={syscall},
    description={Lời gọi hệ thống, giao diện giữa tiến trình người dùng và nhân hệ điều hành},
    first={lời gọi hệ thống (system call - syscall)}
}

\newglossaryentry{ebpf}{
    type=\acronymtype,
    name={eBPF},
    description={extended Berkeley Packet Filter, cơ chế chạy chương trình kiểm tra an toàn trong nhân Linux},
    first={extended Berkeley Packet Filter (eBPF)}
}

\newglossaryentry{bpf}{
    type=\acronymtype,
    name={BPF},
    description={Berkeley Packet Filter, nền tảng gốc của eBPF},
    first={Berkeley Packet Filter (BPF)}
}

\newglossaryentry{dfa}{
    type=\acronymtype,
    name={DFA},
    description={Automata hữu hạn đơn định (Deterministic Finite Automaton)},
    first={automata hữu hạn đơn định (Deterministic Finite Automaton - DFA)}
}

\newglossaryentry{stide}{
    type=\acronymtype,
    name={STIDE},
    description={Sequence Time-Delay Embedding, phương pháp phát hiện bất thường dựa trên chuỗi con syscall},
    first={Sequence Time-Delay Embedding (STIDE)}
}

\newglossaryentry{rnn}{
    type=\acronymtype,
    name={RNN},
    description={Mạng nơ-ron hồi quy (Recurrent Neural Network)},
    first={mạng nơ-ron hồi quy (Recurrent Neural Network - RNN)}
}

\newglossaryentry{lstm}{
    type=\acronymtype,
    name={LSTM},
    description={Long Short-Term Memory, một biến thể của mạng nơ-ron hồi quy},
    first={Long Short-Term Memory (LSTM)}
}

\newglossaryentry{gru}{
    type=\acronymtype,
    name={GRU},
    description={Gated Recurrent Unit, một biến thể gọn hơn của mạng nơ-ron hồi quy},
    first={Gated Recurrent Unit (GRU)}
}

\newglossaryentry{xai}{
    type=\acronymtype,
    name={XAI},
    description={Trí tuệ nhân tạo giải thích được (Explainable Artificial Intelligence)},
    first={trí tuệ nhân tạo giải thích được (Explainable Artificial Intelligence - XAI)}
}

\newglossaryentry{auroc}{
    type=\acronymtype,
    name={AUROC},
    description={Diện tích dưới đường cong ROC (Area Under the Receiver Operating Characteristic Curve)},
    first={Area Under the Receiver Operating Characteristic Curve (AUROC)}
}

\newglossaryentry{f1}{
    type=\acronymtype,
    name={F1},
    description={Thước đo trung hòa giữa precision và recall},
    first={F1 score}
}

\newglossaryentry{tid}{
    type=\acronymtype,
    name={TID},
    description={Định danh luồng thực thi (Thread ID)},
    first={định danh luồng (Thread ID - TID)}
}

\newglossaryentry{cpu}{
    type=\acronymtype,
    name={CPU},
    description={Bộ xử lý trung tâm (Central Processing Unit)},
    first={Central Processing Unit (CPU)}
}

\newglossaryentry{gpu}{
    type=\acronymtype,
    name={GPU},
    description={Bộ xử lý đồ họa, thường dùng để tăng tốc huấn luyện mô hình học sâu (Graphics Processing Unit)},
    first={Graphics Processing Unit (GPU)}
}

\newglossaryentry{mitre-attack}{
    type=\acronymtype,
    name={MITRE ATT\&CK},
    description={Khung tri thức mô tả chiến thuật và kỹ thuật tấn công mạng trong thực tế},
    first={MITRE ATT\&CK}
}
